\documentclass[UTF8,a4paper,10.5pt]{ctexart}
\usepackage[left=1.9cm,right=1.9cm,top=1.9cm,bottom=2.0cm]{geometry}
\usepackage{amsmath,amssymb}
\usepackage{enumitem}
\usepackage{fancyhdr}
\usepackage{lastpage}
\usepackage{listings}
\usepackage{xcolor}
\usepackage{array}

\pagestyle{fancy}
\fancyhf{}
\renewcommand{\headrulewidth}{0pt}
\fancyfoot[C]{第~\thepage~页\quad 共~\pageref{LastPage}~页}

\lstset{
    language=C,
    basicstyle=\ttfamily\small,
    keywordstyle=\color{blue!70!black}\bfseries,
    commentstyle=\color{green!45!black},
    stringstyle=\color{red!60!black},
    columns=fullflexible,
    keepspaces=true,
    breaklines=true,
    showstringspaces=false,
    escapeinside={(*@}{@*)},
    aboveskip=3pt,
    belowskip=3pt
}

\newcommand{\blank}[1][2.8cm]{\underline{\hspace{#1}}}
\newcommand{\scorebox}[1]{%
    \noindent
    \begin{tabular}{|c|c|}
        \hline
        分\quad 数 & \hspace{1.1cm} \\
        \hline
        评卷人 & \hspace{1.1cm} \\
        \hline
    \end{tabular}%
    \hspace{0.8em}%
    \parbox[c]{\dimexpr\linewidth-4.4cm\relax}{\textbf{#1}}%
    \par\medskip
}
\newcommand{\parenans}{\hfill(\hspace{1.4cm})}

\begin{document}

\begin{center}
    {\large\bfseries 华中科技大学人工智能与自动化学院 2020\textasciitilde 2021 学年第一学期}\\[4pt]
    {\Large\bfseries “C语言程序设计”考试试卷（A 卷）}\\[8pt]
    考试方式：\underline{\hspace{0.6em}闭卷\hspace{0.6em}}\qquad
    考试日期：\underline{\hspace{0.6em}2021-01-08\hspace{0.6em}}\qquad
    考试时长：\underline{\hspace{0.5em}150\hspace{0.5em}} 分钟\\[8pt]
    \begin{tabular}{|*{8}{>{\centering\arraybackslash}p{1.4cm}|}}
        \hline
        题号 & 一 & 二 & 三 & 四 & 五 & 六 & 总分 \\
        \hline
        分数 & & & & & & & \\[4pt]
        \hline
    \end{tabular}
\end{center}

\bigskip

\scorebox{一、判断下列语句或程序段的对错。（“$\times$”表示错，“$\checkmark$”表示对）（10分）}

\begin{enumerate}[label=(\arabic*), leftmargin=2.2em, itemsep=2pt]
    \item \texttt{int x=05678;} \parenans
    \item \texttt{int x = 0x177f;}\\
          \texttt{float y = 1.8e+5;} \parenans
    \item \lstinline|char _a01='\101';| \parenans
    \item \texttt{int **p, (*pa)[10];}\\
          \texttt{p = pa;} \parenans
    \item \texttt{int (*pa)[5],a[5][5];}\\
          \texttt{pa = a++;} \parenans
    \item \texttt{char *str[10];}\\
          \texttt{gets(str[0]);} \parenans
    \item \texttt{int *p, a[5][5];}\\
          \texttt{p = *a+1;} \parenans
    \item \texttt{int a[5]=\{0, 1, 2,3,4,5\};} \parenans
    \item \texttt{char **p, *pa[10];}\\
          \texttt{p = pa;} \parenans
    \item \lstinline|char *str = "\\tabcd\123\\n";| \parenans
\end{enumerate}

\bigskip

\scorebox{二、计算下列表达式的值（10分）}

\noindent 假设 \texttt{int} 和 \texttt{unsigned int} 类型均为 16 位长度，且各题彼此独立。

\noindent 设 \texttt{unsigned int a=2, b=3; int c=4, d=5; float f; char str[10]="hust";}
\begin{enumerate}[label=(\arabic*), leftmargin=2.2em, itemsep=2pt]
    \item \texttt{f = c/b + d/a;} \parenans
    \item \texttt{--a || --a \&\& a;} \parenans
    \item \texttt{(a\^{}c) \& (b|d)} \parenans
    \item \texttt{f = sizeof(str) + strlen(str);} \parenans
    \item \texttt{a-c > d - b ? a : b;} \parenans
\end{enumerate}

\bigskip

\scorebox{三、程序改错（10分）要求：不得改变程序框架，不得重写程序，无需文字说明，直接在代码上添加、删除和修改。}

\noindent 1、如下程序对已知的 \texttt{N} 个字符串排序，其中在 \texttt{main()} 函数中初始化，在 \texttt{sort} 函数中排序（使用指针数组）。
\begin{lstlisting}
#include <stdio.h>

void main()
{
    int i;
    char str[N][80] = {"basic", "pascal", "colbol", "losp", "agol"};
    char *pstr[N];
    sort(pstr);
    for (i = 0; i < N; i++)
    {
        printf("%s\n", str[i]);
    }
}

void sort(char (*p)[80])
{
    int i, j;
    char *temp;
    for (i = 0; i < N; i++)
        for (j = 0; j < N - 1; j++)
            if (p[j] > p[j + 1])
            {
                temp = p[j];
                p[j] = p[j + 1];
                p[j + 1] = temp;
            }
}
\end{lstlisting}

\medskip
\noindent 2、设有 4 个候选人，\texttt{N} 个人参加选举，每次输入一个得票的候选人的名字，要求最后输出个人的得票结果。
\begin{lstlisting}
#include <stdio.h>

#define N 10
struct person
{
    char name;
    int count;
};

void main()
{
    struct person leader[4] = {{"wang", 0},
                               {"zhao", 0},
                               {"mao",  0},
                               {"gao",  0}};
    int i, j;
    for (i = 0; i < N; i++)
    {
        gets(name);
        for (j = 0; j < 4; j++)
            if (name == leader[i].name)
            {
                count++;
                continue;
            }
    }
    printf("\n");
    for (j = 0; j < 4; j++)
        printf("%s:%d\n", leader[j].name, leader[j].count);
}
\end{lstlisting}

\bigskip

\scorebox{四、程序填空（10分）}

\noindent (1)~求已知两个正整数的最大公约数。
\begin{lstlisting}
#include <stdio.h>

void main()
{
    (*@\blank[3.2cm]@*);
    int a, b;

    printf("please input a ,b:");
    scanf("%d%d", (*@\blank[3.2cm]@*));

    for (i = (*@\blank[3.2cm]@*); i > 0; i--)  // i 初值为 a,b 中的较小值
    {
        if ((*@\blank[3.8cm]@*))
        {
            printf("the max  = %d ", i);
            (*@\blank[3.2cm]@*);
        }
    }
}
\end{lstlisting}

\medskip
\noindent (2)~求序列：$1!,\ 2!,\ 3!,\ 4!\cdots\cdots$ 的前 10 项之和。
\begin{lstlisting}
(*@\blank[3.5cm]@*)

void main()
{
    int i;
    long int sum, t;  // sum 代表和, t 代表某项

    sum = 0;
    t = 1;

    for (i = 1; (*@\blank[3.2cm]@*); i++)
    {
        sum = (*@\blank[3.2cm]@*);
        (*@\blank[4.2cm]@*);
    }
    printf(" sum = %ld", (*@\blank[3.2cm]@*));
}
\end{lstlisting}

\bigskip

\scorebox{五、输出程序运行结果（25分，每题5分，结果写在题目的右边）}

\noindent 1、
\begin{lstlisting}
#include <stdio.h>

void main()
{
    int i;
    for (i = 1; i < 10; i++)
    {
        if (i % 3)
        {
            putchar('W');
        }
        else if (i % 6 == 0) break;
        else continue;
        printf("U  HAN\n");
    }
    printf("I  love you");
}
\end{lstlisting}

\medskip
\noindent 2、
\begin{lstlisting}
#include <stdio.h>

int n;

void display()
{
    int a = 1;
    static int b = 1;
    printf("a=%d,  b=%d  c=%d\n", a++, b++, ++n);
}

void main()
{
    int i;
    for (i = 0; i < 5; i++)
    {
        display();
    }
}
\end{lstlisting}

\newpage
\noindent 3、
\begin{lstlisting}
#include <stdio.h>

void main()
{
    char str[][20] = {"china",
                      "hubei",
                      "wuhan",
                      "hust-aia",
                      "c language"};
    int i;
    for (i = 0; i < 5; i++)
    {
        printf("%s\n", str[4 - i] + 4 - i);
    }
}
\end{lstlisting}

\medskip
\noindent 4、
\begin{lstlisting}
#include <stdio.h>

void main()
{
    char *str[] = {"pascal", "cobol", "fortran", "lisp", "java"};
    char **p[] = {str + 4, str + 3, str + 2, str + 1, str};
    char ***pp = p;

    printf("%s---", **(pp++) + 1);
    printf("%s---", **(pp++) + 2);
    printf("%s---", **(pp++) + 1);
    printf("%s",    **(pp++) + 2);
    printf("---%s", **(pp));
}
\end{lstlisting}

\medskip
\noindent 5、
\begin{lstlisting}
#include <stdio.h>
#define N 6

void main()
{
    int a[N][N];
    int i, j;

    for (i = 0; i < N; i++)
    {
        a[i][0] = 1;
        a[i][i] = 1;
    }
    for (i = 2; i < N; i++)
    {
        for (j = 1; j < i; j++)
        {
            a[i][j] = a[i - 1][j - 1] + a[i - 1][j];
        }
    }
    for (i = 0; i < N; i++)
    {
        for (j = 0; j <= i; j++)
        {
            printf("%5d", a[i][j]);
        }
        printf("\n");
    }
}
\end{lstlisting}

\bigskip
\scorebox{六、编程（35分）注意：不得使用全局变量，注意程序结构}

\noindent 1、(1) 编一程序，输入 $x$ 的值，输出 $y$ 的值。（6分）
\[
y = \begin{cases}
x + 10, & x < 10 \\
2x + 200, & 10 \le x < 100 \\
3x + 300, & x \ge 100
\end{cases}
\]

\medskip
\noindent 2、编写一猜数游戏程序，随机产生某个小于 $N$ 的正整数，从键盘反复输入整数进行猜数，当未猜中时，提示输入过大或过小。猜中时，指出猜的次数。最多允许猜 20 次。（8分）

\medskip
\noindent 3、编一程序，要求主函数中输入一行英文和一个字母，被调用的函数返回删除该字母后的那行英文（如有多个，一并删除）并在主函数中输出。（9分）

\medskip
\noindent 4、某班有学生 $N$ 名，每名学生信息由姓名、性别、学号、成绩 1、成绩 2 和成绩 3 组成，试编程要求：
\begin{enumerate}[label=（\arabic*）, leftmargin=2.4em, itemsep=2pt]
    \item 学生信息由键盘输入
    \item 求每人平均分
    \item 按平均分数从高到低排序
    \item 按平均分从高到低显示所有平均分不及格的男生信息。
\end{enumerate}
\noindent （每小题写一函数，通过 \texttt{main} 函数调用实现）（12分）

\end{document}
